\section{Limitations: what these studies do and do not establish}
\label{sec:limitations}

These controlled tasks demonstrate context-dependent predictive relationships,
not human-like understanding or a general world model. Real settings contain
strategic interaction, unreliable cues, changing variables, and longer feedback
horizons absent from our market packets and constructed domains.

The internal intervention concerns one checkpoint and one convention; the
historical study covers one market cohort and training recipe. Additional
architectures, later cohorts, interactive settings, and misleading evidence
would test generality.

Performance is heterogeneous. Smaller checkpoints often struggle to induce the
KIV/ZOR meaning. The episode-matching advantage under joint domain and mechanism
shift does not survive the eight-seed replication; the benefit for compositions
of familiar mechanisms is checkpoint-specific and uses a fixed prior, not a
distribution-matched shuffled-target control. Structure diagnostics use no target
observations; the supervised pilot has one seed and differs from RL in update
budget and, at the successful setting, learning rate. Semantic cue-to-pattern
learning could explain its in-domain success. The gain intervention uses only
two retained checkpoint seeds and does not establish held-out evidence use.
Historical-market gains narrow at the strongest Qwen base. The nonmonotonic
scale comparison uses one recipe; larger models may require different
optimization or data budgets.
